\section{Introduction}
\label{sec:intro}

The Networking and Cryptography library (NaCl)~\cite{BLS12}
is an easy-to-use, high-security, high-speed cryptography library.
It uses specialized code for different platforms, which makes it rather complex and hard to audit.
TweetNaCl~\cite{BGJ+15} is a compact re-implementation in C
of the core functionalities of NaCl and is claimed to be
\emph{``the first cryptographic library that allows correct functionality
to be verified by auditors with reasonable effort''}~\cite{BGJ+15}.
The original paper presenting TweetNaCl describes some effort to support
this claim, for example, formal verification of memory safety, but does not actually
prove correctness of any of the primitives implemented by the library.

One core component of TweetNaCl (and NaCl) is the key-exchange protocol X25519 presented
by Bernstein in~\cite{rfc7748}.
This protocol is standardized in RFC~7748 and used by a wide variety of applications~\cite{things-that-use-curve25519}
such as SSH, Signal Protocol, Tor, Zcash, and TLS to establish a shared secret over
an insecure channel.
The X25519 key-exchange protocol is an \xcoord-only
elliptic-curve Diffie--Hellman key exchange using the Montgomery
curve $E: y^2 = x^3 + 486662 x^2 + x$ over the field $\F{\p}$.
Note that originally, the name ``Curve25519'' referred to the key-exchange protocol,
but Bernstein suggested to rename the protocol to X25519 and to use the name
Curve25519 for the underlying elliptic curve~\cite{Ber14}.
We use this updated terminology in this paper.

\subheading{Contribution of this paper.}
We provide a mechanized formal proof of the correctness of the X25519
implementation in TweetNaCl.
This proof is done in three steps:

We first formalize RFC~7748~\cite{rfc7748} in Coq~\cite{coq-faq}.
% This part uses generic techniques to facilitate later proofs.

In the second step we prove equivalence of the C implementation of X25519
to our RFC formalization.
This part of the proof uses the Verifiable Software Toolchain (VST)~\cite{2012-Appel}
to establish a link between C and Coq.
VST uses Separation logic~\cite{1969-Hoare,Reynolds02separationlogic}
to show that the semantics of the program satisfy a functional specification in Coq.
To the best of our knowledge, this is the first time that
VST is used in the formal proof of correctness of an implementation
of an asymmetric cryptographic primitive.

In the last step we prove that the Coq implementation matches
the mathematical definition of X25519 as given in~\cite[Sec.~2]{Ber06}.
This provides additional confidence that there are no transcription errors
in our formalized version of X25519.
We do this by extending the Coq library
for elliptic curves~\cite{BartziaS14} by Bartzia and Strub to
support Montgomery curves, and in particular Curve25519.
This extension may be of independent interest.

\subheading{Related work.}
Two methodologies exist to get formal guarantees that software meets its specification.
The first is to write a formal specification and then synthesize
machine code by refinement; the second is to formalize
a specification and then verify that some
implementation satisfies it.

The synthesis approach was used by Zinzindohou{\'{e}}, Bartzia, and Bhargavan to build a verified extensible
library of elliptic curves~\cite{Zinzindohoue2016AVE} in F*~\cite{DBLP:journals/corr/BhargavanDFHPRR17}.
This served as ground work for the cryptographic library HACL*~\cite{zinzindohoue2017hacl}
used in the NSS suite from Mozilla.

Coq not only allows verification but also synthesis~\cite{CpdtJFR}.
Erbsen, Philipoom, Gross, and Chlipala make use of it to have
correct-by-construction finite field arithmetic, which is used
to synthesize certified elliptic-curve crypto software~\cite{Philipoom2018CorrectbyconstructionFF,Erbsen2017CraftingCE,Erbsen2016SystematicSO}.
This software suite is now being used in BoringSSL~\cite{fiat-crypto}.

The verification approach has been used to prove the correctness of OpenSSL's
implementations of HMAC~\cite{Beringer2015VerifiedCA}, mbedTLS HMAC-DRBG~\cite{2017-Ye} and SHA-256~\cite{2015-Appel}.
In terms of languages and tooling, this work is closest to what we present here,
but our work considers an asymmetric primitive and provides computer-verified
proofs up to the mathematical definition of the group theory behind elliptic curves.

When considering the target of the verification effort, the closest work
to this paper is presented in~\cite{Chen2014VerifyingCS}, which presents
a mechanized proof of two assembly-level implementations of the core function of X25519.
The proof in~\cite{Chen2014VerifyingCS} takes a different approach from ours.
It uses heavy annotation of code together with SAT solvers;
also, it does not cover the full X25519 functionality and does
not make the link to the mathematical definition from~\cite{Ber06}.

%Synthesis approaches force the user to depend on generated code which may not
%be optimal (speed, compactness) and they require compiler tweaks in order
%to achieve the desired properties. On the other hand, carefully handcrafted
%optimized code forces the verifier to consider all the possible special cases
%and write a low-level specification of the details of the code in order to prove
%the correctness of the algorithm.

%It has to be noted that Curve25519 implementations have already been under scrutiny
%However in this paper we provide a proofs of correctness and soundness of a C program up to
%the theory of elliptic curves.

\subheading{Reproducing the proofs.}
To maximize reusability of our results we placed the code of our formal proof
presented in this paper into the public domain.
It is available at \url{https://cdn-09.anonfile.com/6fp8ta70n9/014c4f8c-1569921449/coq-verif-tweetnacl.tar.gz}
with instructions of how to compile and verify our proof.
A description of the content of the code archive is provided in
Appendix~\ref{appendix:proof-folders}.

\subheading{Organization of this paper.}
\sref{sec:preliminaries} gives the necessary background on Curve25519 and X25519
implementations and a brief explanation of how formal verification works.
\sref{sec:Coq-RFC} provides our Coq formalization of X25519 as specified in RFC~7748~\cite{rfc7748}.
\sref{sec:C-Coq} provides the specifications of X25519 in TweetNaCl and some of the
proof techniques used to show the correctness with respect to RFC~7748~\cite{rfc7748}.
\sref{sec:maths} describes our extension of the formal library by Bartzia
and Strub and the correctness of X25519 implementation with respect to Bernstein's
specifications~\cite{Ber14}.
Finally in \sref{sec:Conclusion} we discuss the trusted code base of our proofs.
\fref{tikz:ProofOverview} shows a graph of dependencies of the proofs.

\begin{figure}[h]
  \centering
  \begin{tikzpicture}[textstyle/.style={black, anchor= south west, align=center, minimum height=0.45cm, text centered, font=\scriptsize}]


  % adapt line thickness and line width, if needed

   \path [thick, dashed] (2.5,1) edge +(0,-6.75);
   \draw (2.5,1) node[anchor=north east] {\sref{sec:Coq-RFC}};
   \draw (2.5,1) node[anchor=north west] {\sref{sec:C-Coq}};
   \path [thick, dashed] (0,-5.75) edge +(8.5,0);
   \draw (8.5,-5.75) node[anchor=north east] {\sref{sec:maths}};

  % SECTION III

  % Definition of RFC
  \begin{scope}[yshift=-3 cm,xshift=0 cm]
   \draw (0,0) -- (0.4, 0.4) -- (1.75,0.4) -- (1.75,0) -- cycle;
   \draw (0,0) -- (1.75,0) -- (1.75,-1) -- (0, -1) -- cycle;
   \draw (0.3,-0.05) node[textstyle] {Definition};
   \draw (0.875,-0.5) node[textstyle, anchor=center] {\texttt{RFC}};
   \draw (1.75,-1) node[textstyle, anchor = south east] {\color{doc@lstfunctions}{.V}};
  \end{scope}

   % SECTION IV

  % C code
  \begin{scope}[yshift=-0.25 cm,xshift=2.75 cm]
    \draw (0,0) -- (0.4, 0.4) -- (1.25,0.4) -- (1.25,0) -- cycle;
    \draw (0,0) -- (1.25,0) -- (1.25,-1.25) -- (0, -1.25) -- cycle;
    \draw (0.3,-0.05) node[textstyle] {code};
    \draw (0.675,-0.5) node[textstyle, anchor=center] {\texttt{Prog}};
    \draw (1.25,-1.25) node[textstyle, anchor = south east] {\color{doc@lstfunctions}{.C}};
  \end{scope}

  % V code
  \begin{scope}[yshift=-0.25 cm,xshift=5 cm]
    \draw (0,0) -- (0.4, 0.4) -- (1.25,0.4) -- (1.25,0) -- cycle;
    \draw (0,0) -- (1.25,0) -- (1.25,-1.25) -- (0, -1.25) -- cycle;
    \draw (0.3,-0.05) node[textstyle] {code};
    \draw (0.675,-0.5) node[textstyle, anchor=center] {\texttt{Prog}};
    \draw (1.25,-1.25) node[textstyle, anchor = south east] {\color{doc@lstfunctions}{.V}};
    % \draw (1.25,0)  node[anchor=south east, inner sep=0pt] {\includegraphics[width=.0125\textwidth]{img/coq_logo.png}};
  \end{scope}

  \path [thick, dotted, double, ->] (4,-0.75) edge (5, -0.75);

  % VST Spec
  \begin{scope}[yshift=-3 cm,xshift=2.75 cm]
    \draw (0,0) -- (0.4, 0.4) -- (2.75,0.4) -- (2.75,0) -- cycle;
    \draw (0,0) -- (2.75,0) -- (2.75,-2) -- (0, -2) -- cycle;
    \draw (0.3,-0.05) node[textstyle] {Specification};
    \draw (1.375,-1) node[textstyle, anchor=center, align=left] {
      {\color{doc@lstnumbers}\textbf{Pre}}:\\
      ~~\texttt{s}$[\!\!\{$\texttt{n}$\}\!\!]\leftarrow n \in \N$,\\
      % ~~$n \in$ \TNaCles{u8[32]},\\
      ~~\texttt{s}$[\!\!\{$\texttt{p}$\}\!\!]\leftarrow P \in E(\F{p^2})$\\
      % ~~$P \in$ \TNaCles{u8[32]}\\
      {\color{doc@lstdirective}\textbf{Post}}:\\
      ~~\texttt{s}$[\!\!\{$\texttt{q}$\}\!\!]\leftarrow$~\texttt{RFC}$(n,P)$};
      \draw (2.75,-2) node[textstyle, anchor = south east] {\color{doc@lstfunctions}{.V}};
  \end{scope}

  % VST Theorem
  \begin{scope}[yshift=-3 cm,xshift=6 cm]
    \draw (0,0) -- (0.4, 0.4) -- (2.5,0.4) -- (2.5,0) -- cycle;
    \draw (0,0) -- (2.5,0) -- (2.5,-1.25) -- (0, -1.25) -- cycle;
    \draw (0.3,-0.05) node[textstyle] {Proof};
    \draw (1.25,-0.5) node[textstyle, anchor=center] {\{{\color{doc@lstnumbers}\textbf{Pre}}\} \texttt{Prog} \{{\color{doc@lstdirective}\textbf{Post}}\}};
    \draw (2.5,0) node[textstyle, anchor = south east] {\color{doc@lstfunctions}{\checkmark}};
    \draw (2.5,-1.25) node[textstyle, anchor = south east] {\color{doc@lstfunctions}{.V}};
  \end{scope}

  \path [thick, double] (5.625,-1.5) edge [out=-90, in=90] (5.625, -2.5);
  \path [thick, double, ->] (5.625, -2.5) edge [out=-90, in=180] (6, -3.5);
  \path [thick, double, ->] (5.5,-3.75) edge [out=0, in=180] (6, -3.75);



  % SECTION V

  % Spec of Curve nP
  \begin{scope}[yshift=-7.5 cm,xshift=0 cm]
    \draw (0,0) -- (0.4, 0.4) -- (1.75,0.4) -- (1.75,0) -- cycle;
    \draw (0,0) -- (1.75,0) -- (1.75,-1) -- (0, -1) -- cycle;
    \draw (0.3,-0.05) node[textstyle] {Definition};
    \draw (0.875,-0.5) node[textstyle, anchor=center] {$n \cdot P$};
    \draw (1.75,-1) node[textstyle, anchor = south east] {\color{doc@lstfunctions}{.V}};
  \end{scope}

  % Correctness Theorem
  \begin{scope}[yshift=-7 cm,xshift=6 cm]
    \draw (0,0) -- (0.4, 0.4) -- (2.5,0.4) -- (2.5,0) -- cycle;
    \draw (0,0) -- (2.5,0) -- (2.5,-1.5) -- (0, -1.5) -- cycle;
    \draw (0.3,-0.05) node[textstyle] {Proof};
    \draw (1.25,-0.75) node[textstyle, anchor=center] {{\color{doc@lstnumbers}\textbf{Pre}} $\implies$\\$\text{\texttt{RFC}}(n,P) = n \cdot P$};
    \draw (2.5,0) node[textstyle, anchor = south east] {\color{doc@lstfunctions}{\checkmark}};
    \draw (2.5,-1.5) node[textstyle, anchor = south east] {\color{doc@lstfunctions}{.V}};
  \end{scope}

  \path [thick, double, ->] (1.75,-3.5) edge  [out=0, in=-180] (2.75, -3.5);
  \path [thick, double] (1.75,-3.5) edge [out=0, in=90] (2.25, -4);
  \path [thick, double] (2.25, -4) edge [out=-90, in=90] (2.25, -7);
  \path [thick, double] (2.25, -7) edge [out=-90, in=-180] (3, -7.75);
  \path [thick, double, ->] (3, -7.75) edge [out=0, in=-180] (6, -7.75);
  \path [thick, double, ->] (1.75,-8) edge [out=0, in=-180] (6, -8);


\end{tikzpicture}

  \caption{Structure of the proof.}
  \label{tikz:ProofOverview}
\end{figure}

\todoln{NEW}
In this figure, {\color{doc@lstfunctions}.V} represent Coq files.

% Five years ago:
% \url{https://www.imperialviolet.org/2014/09/11/moveprovers.html}
% \url{https://www.imperialviolet.org/2014/09/07/provers.html}

% \section{Related Works}
%
% \begin{itemize}
%   \item HACL*
%   \item Proving SHA-256 and HMAC
%   \item \url{http://www.iis.sinica.edu.tw/~bywang/papers/ccs17.pdf}
%   \item \url{http://www.iis.sinica.edu.tw/~bywang/papers/ccs14.pdf}
%   \item \url{https://cryptojedi.org/crypto/#gfverif}
%   \item \url{https://cryptojedi.org/crypto/#verify25519}
%   \item Fiat Crypto : synthesis
% \end{itemize}
%
% Add comparison with Fiat-crypto
